\section{Proof sketches}
\label{sec:proof_sketches}

This section provides proof sketches for the classification results.
All arguments are existence- and invariance-based. No reconstruction,
inference, optimization, algorithmic procedures, or asymptotic limits
are used.

\subsection{Proof sketch for Proposition~\ref{prop:observational-invariance}
(observational invariance)}

Assume that $\varphi = \tilde{\varphi} \circ \Pi$ for some mapping
\[
\tilde{\varphi} : \Sigma_{\mathrm{obs}} \to \{\text{true},\text{false}\}.
\]
Let $K_{EA}^{(1)}$ and $K_{EA}^{(2)}$ be observationally equivalent
enterprise continua, i.e.
\[
\Pi(\Omega_{EA}^{(1)}) = \Pi(\Omega_{EA}^{(2)}).
\]

For any internally admissible state
$s \in \Omega_{EA}^{(1)} \cup \Omega_{EA}^{(2)}$, the observable
signature $\Pi(s)$ belongs to the same observable signature set in both
continua. Since $\tilde{\varphi}$ depends exclusively on the observable
signature, the value of $\varphi(s)$ cannot differ across members of the
same observational equivalence class.

Therefore, $\varphi$ is invariant across all observationally equivalent
enterprise continua and is decidable under incomplete observability.

The argument is purely extensional: decidability follows solely from the
fact that the predicate factors through the observation operator.

\subsection{Proof sketch for Theorem~\ref{thm:stop-decidability}
(decidability of STOP boundary predicates)}

The theorem assumes that STOP membership corresponds to a boundary
condition in $\partial\Omega_{EA}$ whose observable consequences are
invariant across any observational equivalence class.

Under this assumption, if a given observable signature $\sigma \in
\Sigma_{\mathrm{obs}}$ is compatible with STOP for at least one
internally admissible enterprise state, then every internally
admissible state compatible with $\sigma$ must also lie in the STOP
region.

Hence STOP membership becomes a class property of observational
equivalence, rather than a property of individual internal states.
The corresponding predicate $\varphi_{\mathrm{STOP}}$ is therefore
decidable under incomplete observability.

This proof relies only on the ETS-defined relationship between STOP and
$\partial\Omega_{EA}$, together with the stated invariance condition.
No additional structural assumptions are introduced.

\subsection{Proof sketch for Theorem~\ref{thm:structural-undecidability}
(structural undecidability)}

Let $\varphi$ be a predicate that depends on at least one internal
OC/ETS element not uniquely determined by observable signatures.

To establish undecidability, it suffices to exhibit two
observationally equivalent enterprise continua
$K_{EA}^{(1)}$ and $K_{EA}^{(2)}$ such that
\[
\varphi(\Omega_{EA}^{(1)}) \neq \varphi(\Omega_{EA}^{(2)}).
\]

The construction follows a standard indistinguishability argument:

\begin{enumerate}
\item Fix an admissible observable signature set
      $S \subseteq \Sigma_{\mathrm{obs}}$.
\item Consider the observational consistency set
      $\Pi^{-1}(S) \subseteq \Omega_{EA}$ containing all internally
      admissible states compatible with $S$.
\item By hypothesis, $\varphi$ depends on an internal feature not fixed
      by $S$. Therefore, there exist internally admissible states
      $s^{(1)}, s^{(2)} \in \Pi^{-1}(S)$ such that
      \[
      \varphi(s^{(1)}) \neq \varphi(s^{(2)}).
      \]
\item Lift these states into two enterprise continua realizing the same
      observable signature set $S$ but differing with respect to the
      relevant internal feature. Then
      \[
      \Pi(\Omega_{EA}^{(1)}) = \Pi(\Omega_{EA}^{(2)})
      \quad\text{while}\quad
      \varphi(\Omega_{EA}^{(1)}) \neq \varphi(\Omega_{EA}^{(2)}).
      \]
\end{enumerate}

This establishes undecidability directly from the existence of
non-trivial observational equivalence classes.

\subsection{Proof sketch for Corollary~\ref{cor:nonterminal-undecidability}
(non-terminal phase distinctions)}

Assume that two non-STOP phase outcomes (e.g.\ OK and INERTIA) are
defined via internal conditions that are not observationally invariant.

The predicate distinguishing these outcomes therefore depends on
internal features not uniquely determined by
$\Sigma_{\mathrm{obs}}$. By
Theorem~\ref{thm:structural-undecidability}, there exist
observationally equivalent enterprise continua that differ with respect
to the phase predicate.

Consequently, no admissible observation can decide the distinction
between such non-terminal phase outcomes.

\subsection{Proof sketch for Theorem~\ref{thm:stop-only}
(STOP-only classification)}

Let $\varphi$ be a phase predicate taking values in
$\{\text{OK}, \text{INERTIA}, \text{COLLAPSE}, \text{STOP}\}$.
Assume that STOP membership is observationally invariant, while the
remaining phase values are not.

\begin{itemize}
\item \emph{STOP decidability.} By invariance, if any internally
      admissible state compatible with an observable signature is STOP,
      then all internally admissible states compatible with that same
      signature are STOP. Hence STOP is decidable.
\item \emph{Non-STOP undecidability.} For any non-STOP value, lack of
      invariance implies the existence of observationally equivalent
      enterprise continua for which $\varphi$ takes different non-STOP
      values. Hence non-STOP outcomes are undecidable.
\end{itemize}

Therefore, $\varphi$ is STOP-only decidable.

\subsection{Proof sketch for Proposition~\ref{prop:exhaustive-trichotomy}
(exhaustiveness)}

For any admissible predicate $\varphi$, consider its behavior on each
observational equivalence class induced by $\Pi$.

\begin{enumerate}
\item If $\varphi$ is constant on every equivalence class, then it is
      decidable by Definition~\ref{sec:definitions}.
\item If $\varphi$ is non-constant on at least one equivalence class,
      then it is undecidable by Definition~\ref{sec:definitions},
      unless a one-sided invariance holds for STOP.
\item If STOP is invariant while non-STOP values are not, then
      $\varphi$ is STOP-only decidable by
      Definition~\ref{sec:definitions}.
\end{enumerate}

These cases exhaust all possibilities, yielding the stated trichotomy.
